\title{xLSTM: Extended Long Short-Term Memory}

\begin{abstract}
The fundamental concepts behind the Long Short-Term Memory (LSTM), namely the constant error carousel and gating mechanisms, emerged in the 1990s. Over subsequent decades, LSTMs have maintained their relevance, underlining key achievements in deep learning—most notably, serving as the basis for the earliest Large Language Models (LLMs). Nevertheless, the introduction of the Transformer framework, whose hallmark is a highly parallelizable self-attention mechanism, signified a significant shift, enabling scalability beyond what LSTMs previously offered. This work poses a straightforward inquiry: to what extent can language modeling advance by scaling LSTMs to the billion-parameter regime, making use of cutting-edge methods from modern LLMs, while specifically addressing existing weaknesses associated with LSTMs? We propose, first, the application of exponential gating, complemented by refined normalization and stabilization strategies. Second, we adapt the LSTM's memory architecture, yielding: (i) sLSTM, characterized by its scalar memory, scalar updates, and revised memory blending, and (ii) mLSTM, which is compatible with full parallelization and features matrix-based memory with a covariance-style update procedure. These LSTM variants, when incorporated within residual block-based networks, form what we term xLSTM blocks, which can be sequentially stacked to construct xLSTM models. The amalgamation of exponential gating with the enhanced memory architectures significantly elevates the performance and scalability of xLSTM, allowing these models to rival or even surpass contemporary Transformers and State Space Models.\\
Code available at: \url{https://github.com/NX-AI/xlstm}
\end{abstract}

\section{Introduction}
\label{sec:intro}

The Long Short-Term Memory (LSTM) ideas~\citep{Hochreiter:91,Hochreiter:97nips,Hochreiter:97}, i.e., the constant error carousel and gating, were introduced to overcome the vanishing gradient problem of recurrent neural networks~\citep{Hochreiter:91,Hochreiter:00book}:

\begin{align}
\label{eq:lstm_idea}
  c_t \ = \ f_t \ c_{t-1} \ + \ 
          i_t \  z_t \ , \quad  
          h_t \ = \ o_t \ \psi({c_t}) \ . 
\end{align}

In the constant error carousel, the cell state $c_{t-1}$ (indicated in green) receives an additive update through cell inputs $z_t$, with sigmoid gates (in blue) regulating this process. The input gate $i_t$ and forget gate $f_t$ are responsible for managing this update, whereas the output gate $o_t$ determines the memory cell’s output, corresponding to the hidden state $h_t$. The cell state undergoes normalization or is constrained by the function $\psi$, after which the output gate produces the hidden state.

LSTMs have been widely implemented across multiple fields \citep{Hochreiter:01,Hochreiter:07,Schmidhuber:15}, maintaining dominance in text generation tasks until the emergence of Transformers in 2017~\citep{Vaswani:17}. Their robust performance has been established in a variety of sequence-processing applications, including text generation~\citep{Graves:13,Karpathy:15blog}, handwriting synthesis~\citep{Graves:13}, sequence-to-sequence translation~\citep{Sutskever:14nips}, program evaluation~\citep{Zaremba:14arxiv}, image captioning~\citep{Karpathy:15,Hossain:19}, source code generation~\citep{Karpathy:15blog}, rainfall-runoff modeling~\citep{Kratzert:18,Kratzert:19}, and hydrological forecasting for flood warning systems~\citep{Nearing:24}. In the realm of reinforcement learning, LSTMs represent the top-performing sequential architectures, as demonstrated in major systems such as AlphaStar for StarCraft II~\citep{Vinyals:17short}, OpenAI Five for Dota 2~\citep{OpenAI:19}, and control models for nuclear fusion’s magnetic controllers~\citep{Degrave:22short}. LSTMs are especially effective at learning abstract representations, successfully extracting high-level semantic features and retaining them within their memory cells~\citep{Karpathy:15blog}. This capability has been illustrated through the emergence of specialized neurons, such as those sensitive to numerical information and syntax~\citep{Lakretz:19}, linguistic properties~\citep{Bau:19}, and sentiment~\citep{Radford:17arxiv}. Today, LSTMs continue to be employed in impactful applications~\citep{Degrave:22short,Nearing:24} and have proven their longevity and reliability.

Although LSTMs have achieved significant advances, they are subject to three fundamental shortcomings: (i) Inability to update previously stored information. This drawback is illustrated using the {\it Nearest Neighbor Search} task (refer also to Appendix~\ref{sec:appNearestNeighborSearch}): Given a reference vector, the model must scan a sequence to find the vector closest to the reference, ultimately retrieving the value associated with the most similar vector at the end. The left panel in Figure~\ref{fig:lstmProblems} presents the mean squared error for this scenario. Traditional LSTM architectures exhibit difficulty replacing a stored item with a superior candidate found later in the sequence. In contrast, our xLSTM utilizes exponential gating, effectively addressing this shortcoming. (ii) Restricted memory capacity, since LSTM architecture constrains stored information to scalar cell states. This is further demonstrated with the {\it Rare Token Prediction} task. The right panel of Figure~\ref{fig:lstmProblems} reports perplexity results from Wikitext-103~\citep{Merity:17} partitioned by token rarity. The LSTM's perplexity rises on infrequent tokens due to its inherent storage limitations; our xLSTM, employing matrix memory, overcomes these bottlenecks. (iii) Limited ability to parallelize, as information mixing through recurrent hidden-to-hidden pathways mandates inherently sequential computations.

Such limitations have led to the ascent of Transformer models~\citep{Vaswani:17} in the domain of language modeling. If these obstacles are circumvented and LSTMs are extended to match the scale of current Large Language Models, what level of performance can be achieved in language modeling tasks?

\section{Extended Long Short-Term Memory}
\label{sec:xLSTM}

To address the constraints inherent in standard LSTM, the Extended Long Short-Term Memory (xLSTM) model implements two core modifications to the framework depicted in Equation~\eqref{eq:lstm_idea}. Specifically, xLSTM introduces exponential gating and revised memory mechanisms, thereby expanding the LSTM paradigm with two new variants: (i) the sLSTM, as detailed in Section~\ref{sec:sLSTM}, which is characterized by a scalar-valued memory, scalar updates, and an approach based on memory mixing; and (ii) the mLSTM, explained in Section~\ref{sec:mLSTM}, featuring matrix-valued memory and updates performed via the covariance (outer product) operation—a process conducive to full parallelism. Both the sLSTM and mLSTM variants incorporate exponential gating to advance traditional LSTM design. Parallel computation is supported in the mLSTM by eliminating the memory mixing, that is, omitting hidden-to-hidden recurrent interactions. 

Moreover, sLSTM and mLSTM may be generalized to encompass multiple memory cells; the sLSTM applies memory mixing across these cells. In addition, sLSTM is capable of supporting multiple heads—where memory mixing occurs among cells within each head but not between distinct heads—thus establishing an alternative form of memory mixing when used in conjunction with exponential gating. Conversely, for mLSTM, extending the architecture to multiple heads is effectively analogous to introducing multiple cells.

The integration of these variants into residual block modules yields what are termed xLSTM blocks (refer to Section~\ref{sec:xLSTMblock}). Building architectures by stacking xLSTM blocks with residual connections forms the basis of xLSTM architectures, as discussed in Section~\ref{sec:xLSTMarchitecure}. An overview of the xLSTM architecture and its principal components is illustrated in Figure~\ref{fig:xlstm_sketch}.

\subsection{Review of the Long Short-Term Memory}
\label{sec:orgLSTM}

The foundational concept of LSTM~\citep{Hochreiter:91,Hochreiter:97nips,Hochreiter:97} established the scalar memory cell as a core element for both computation and information retention, effectively addressing the issue of vanishing gradients~\citep{Hochreiter:91,Hochreiter:00book} via the constant error carousel mechanism (cell state update). Within this memory cell, three distinct gates regulate its function: the input gate, output gate, and forget gate, the latter being incorporated by \citet{Gers:00}. The dynamics governing the LSTM memory cell at time step $t$ can be described as follows:

\begin{align}
c_t \ &= \  f_t \ c_{t-1} \ + \ i_t \ z_t &  & &\text{cell state} \\
h_t  \ &= \ o_t \ \tilde{h}_t \ , & \tilde{h}_t \ &= \ \psi \left( c_t\right)
  &\text{hidden state} \\
z_t \ &= \ \varphi \left( \tilde{z}_t \right) \ , 
  &\tilde{z}_t \ &=  \ \mathbf{w}^\top_{z} \ \mathbf{x}_t \ + \
  r_{z}  h_{t-1} \ + \  b_{z} \ \
  &\text{cell input} \\
i_t \ &= \ \sigma \left( \tilde{i}_t  \right) \ , 
  &\tilde{i}_t \ &= \ \mathbf{w}^\top_{i} \ \mathbf{x}_t \ + \
  r_{i}  \ h_{t-1} \ + \  b_{i} \ \
  &\text{input gate} \\
f_t \ &= \ \sigma \left(  \tilde{f}_t \right) \ , 
  &\tilde{f}_t \ &= \ \mathbf{w}^\top_{f} \ \mathbf{x}_t  \ + \
  r_{f}  \ h_{t-1} \ + \  b_{f} \ \
  &\text{forget gate} \\
o_t \ &= \ \sigma \left( \tilde{o}_t \right) \ , 
  &\tilde{o}_t  \ &= \ \mathbf{w}^\top_{o} \ \mathbf{x}_t \ + \
  r_{o}  \ h_{t-1} \ + \  b_{o} \ \
  &\text{output gate} 
\end{align}

The vectors $\mathbf{w}_{z}$, $\mathbf{w}_{i}$, $\mathbf{w}_{f}$, and $\mathbf{w}_{o}$ represent the input weights connecting the inputs $\mathbf{x}_t$ to the cell input, input gate, forget gate, and output gate, in that order. Similarly, $r_{z}$, $r_{i}$, $r_{f}$, and $r_{o}$ denote the recurrent weights connecting the previous hidden state $h_{t-1}$ to the cell input, input gate, forget gate, and output gate, respectively. The corresponding bias parameters are given by $b_{z}$, $b_{i}$, $b_{f}$, and $b_{o}$. The functions $\varphi$ and $\psi$ serve as activation functions for the cell input and the hidden state (typically instantiated as $\tanh$). In particular, $\psi$ ensures that the cell state remains within a bounded range; otherwise, the cell state could diverge. Activation functions for all gates are defined as sigmoid functions, i.e., $\sigma \left( x \right)= 1/(1+\exp(-x))$.

Subsequent developments introduced grouping multiple scalar memory cells into a vector $\mathbf{c}_t \in \mathbb{R}^d$, enabling the use of recurrent weight matrices $\mathbf{R} \in \mathbb{R}^{d \times d}$ that allow for interaction among the outputs of different memory cells~\citep{Greff:15}; further elaboration can be found in Appendix~\ref{sec:apporgLSTM}. Investigations through ablation demonstrated that each part of the memory cell is essential~\citep{Greff:15}.

\subsection{sLSTM}
\label{sec:sLSTM}

To enable LSTMs to modify their storage choices, we propose integrating exponential gates (highlighted in red) in combination with normalization and stabilization mechanisms. Specifically, we allow the input and forget gates to utilize exponential activation functions. For normalization purposes, a normalizer state is added, which aggregates the product of the input gate and every subsequent forget gate across future time steps.

The scalar sLSTM forward pass is:

\begin{align}
\label{eq:slstmforward}
c_t \ &= \  f_t \ c_{t-1} \ + \ i_t \ z_t & & 
 &\text{cell state} \\
n_t \ &= \  f_t \ n_{t-1} \ + \ i_t  & & 
 &\text{normalizer state} \\
h_t \ &= \ o_t \ \tilde{h}_t \ ,
  &\tilde{h}_t \ &= \ c_t / n_t
&\text{hidden state} \\
z_t \ &= \ \varphi \left( \tilde{z}_t \right) \ , 
  &\tilde{z}_t \ &=  \ \mathbf{w}^\top_{z} \ \mathbf{x}_t \ + \
  r_{z}  h_{t-1} \ + \  b_{z}
  &\text{cell input} \\
i_t \ &= \ \!\exp \left( \tilde{i}_t  \right) \ , 
  &\tilde{i}_t \ &= \ \mathbf{w}^\top_{i} \ \mathbf{x}_t \ + \
  r_{i}  \ h_{t-1} \ + \  b_{i}
  &\text{input gate} \\
f_t \ &= \ \sigma \left(  \tilde{f}_t \right) \ \text{OR} \
\!\exp \left(  \tilde{f}_t \right) \ ,
  &\tilde{f}_t \ &= \ \mathbf{w}^\top_{f} \ \mathbf{x}_t  \ + \
  r_{f}  \ h_{t-1} \ + \  b_{f}
  &\text{forget gate} \\
o_t \ &= \ \sigma \left( \tilde{o}_t \right) \ , 
  &\tilde{o}_t  \ &= \ \mathbf{w}^\top_{o} \ \mathbf{x}_t \ + \
  r_{o}  \ h_{t-1} \ + \  b_{o}
  &\text{output gate} 
\end{align}

We transfer the original LSTM gating techniques, i.e., input- and/or hidden-dependent gating plus bias term, to the new architectures. Exponential activation functions can lead to large values that cause overflows. Therefore, we stabilize gates with an additional state \mbox{$m_t$~\citep{Milakov:18arxiv}}:

\begin{align}
\label{eq:slstmstabil}
m_t \ &= \ \max \left( \log ( f_t ) + m_{t-1} , \log ( i_t ) \right) &\text{stabilizer state} \\
i'_t \ &= \ \!\exp \left( \log \left ( i_t \right) - m_t \right) \ = \ \!\exp \left( \tilde{i}_t - m_t \right) \ \, 
  &\text{stabil. input gate} \\
f'_t \ &= \ \!\exp \left( \log \left( f_t \right) + m_{t-1} - m_t \right) \ \, 
  &\text{stabil. forget gate}
\end{align}

Appendix~\ref{sec:appsLSTM} demonstrates that substituting $f_t$ with $f'_t$ and $i_t$ with $i'_t$ during the forward pass leaves both the overall network output and the gradients of the loss with respect to the model parameters unchanged.

\paragraph{New Memory Mixing.} Similar to traditional LSTM architectures, sLSTM can accommodate several memory cells, as described in Appendix~\ref{sec:appsLSTM}. Utilizing multiple memory cells facilitates the mixing of memories through recurrent links $\mathbf{R}_{\mathbf{z}}$, $\mathbf{R}_{\mathbf{i}}$, $\mathbf{R}_{\mathbf{f}}$, $\mathbf{R}_{\mathbf{o}}$, which connect the hidden state $\mathbf{h}$ to the memory cell input $\mathbf{z}$ as well as the gates $\mathbf{i}$, $\mathbf{f}$, $\mathbf{o}$. An important novel feature in this mechanism arises from the incorporation of exponential gating. The sLSTM allows for multiple heads, each capable of internal memory mixing but without inter-head mixing. The combination of exponential gating and the multi-head architecture introduces an innovative form of memory mixing unique to sLSTM.

\subsection{mLSTM}
\label{sec:mLSTM}

In order to bolster the memory capacity of LSTMs, we extend the traditional memory cell, which is represented as a scalar $c \in \mathbb{R}$, to a matrix $\mathbf{C} \in \mathbb{R}^{d \times d}$. Consequently, the mechanism for retrieval becomes matrix multiplication. At each time step $t$, our objective is to encode a pair of vectors—specifically, the key $\mathbf{k}_t \in \mathbb{R}^d$ and the value $\mathbf{v}_t \in \mathbb{R}^d$—employing terminology from the Transformer architecture. Subsequently, at a future step $t + \tau$, the original value $\mathbf{v}_t$ is intended to be accessed using a query vector $\mathbf{q}_{t+\tau} \in \mathbb{R}^d$. This protocol is fundamentally aligned with the framework of Bidirectional Associative Memories (BAMs) as described in \citep{Kohonen:72,Anderson:72,Nakano:72,Anderson:77}.

The covariance update rule~\citep{Sejnowski:77,Dayan:91} for storing a key-value pair is

\begin{align}
\mathbf{C}_t \ &= \ \mathbf{C}_{t-1} \ + \ \mathbf{v}_{t} \ \mathbf{k}_t^\top \ .
\end{align}

We consider the application of layer normalization prior to projecting the inputs into key and value vectors, which ensures these vectors are zero-centered. The covariance update mechanism is provably optimal~\citep{Dayan:91} for enhancing the separability among retrieved binary vectors, an objective that aligns with maximizing the signal-to-noise ratio. While even greater separability can be attained by restricting retrieval to pairwise interactions—accepting the associated quadratic complexity characteristic of attention mechanisms~\citep{Krotov:16,Krotov:17,Ramsauer:21}—the covariance update procedure remains fundamentally equivalent to that used by Fast Weight Programmers~\citep{Schmidhuber:92ncfastweights,Schlag:21}. More recent extensions augment this rule with a fixed decay term applied to $\mathbf{C}_{t-1}$, as well as a constant learning rate scaling for 
$\mathbf{v}_{t} \mathbf{k}_t ^\top$~\citep{Ba:16}.
Adopting this perspective, we embed the covariance update formula within the LSTM architecture, such that the forget gate modulates the decay rate, the input gate adjusts the learning rate, and the output gate modulates the magnitude of the retrieved vector.

For this matrix memory, the normalizer state is the weighted sum of key vectors, where each key vector is weighted by the input gate and all future forget gates. Again, the normalizer state keeps record of the strength of the gates. Since the dot product between query and normalizer state can be close to zero, we use the absolute value of this dot product and lower bound it by a threshold (typically 1.0) as done previously~\citep{Sun:23arxiv}. The mLSTM forward pass is:

\begin{align}
\label{eq:mlstm_recurrent_begin}
\mathbf{C}_t \ &= \  f_t \ \mathbf{C}_{t-1} \ + \ 
  i_t \ \mathbf{v}_t \ \mathbf{k}_t^\top & &  &\text{cell state} \\
\mathbf{n}_t \ &= \  f_t \ \mathbf{n}_{t-1} \ + \ 
  i_t \ \mathbf{k}_t & &  &\text{normalizer state} \\
\mathbf{h}_t  \ &= \ \mathbf{o}_t \ \odot \ \tilde{\mathbf{h}}_t
\ , \qquad \qquad 
\tilde{\mathbf{h}}_t \ = \   \mathbf{C}_t \mathbf{q}_t \ / \ 
  \max \left\{ |\mathbf{n}_t^\top \mathbf{q}_t|, 1 \right\} 
   &&&\text{hidden state} \\
\mathbf{q}_t \ &= \ \mathbf{W}_q \ \mathbf{x}_t \ + \ \mathbf{b}_q  & &  &\text{query input} \\
\mathbf{k}_t \ &= \ \frac{1}{\sqrt{d}} \mathbf{W}_k \ \mathbf{x}_t \ + \ \mathbf{b}_k  & &  &\text{key input} \\
\mathbf{v}_t \ &= \ \mathbf{W}_v \ \mathbf{x}_t \ + \ \mathbf{b}_v  & &  &\text{value input} \\
i_t \ &= \ \!\exp \!  \! \left( \tilde{i}_t  \right) \ , 
 \qquad \qquad \ \ \ \ \,
  \tilde{i}_t \ = \ \mathbf{w}^\top_{i} \ \mathbf{x}_t \ + \  b_{i} \ \ \ 
  &&&\text{input gate} \\
f_t \ &= \ \sigma \!  \left(  \tilde{f}_t \right) \ \text{OR} \ \!\exp \!  \! \left(  \tilde{f}_t \right) , \ \ \: \!
\tilde{f}_t \ = \ \mathbf{w}^\top_{f} \ \mathbf{x}_t  \ + \
  b_{f}
  &&& \text{forget gate} \\
\label{eq:mlstm_recurrent_end}
\mathbf{o}_t \ &= \ \sigma \left( \tilde{\mathbf{o}}_t \right) \ , \qquad \qquad \quad \ \ \ \,
  \tilde{\mathbf{o}}_t  \ = \ \mathbf{W}_{\mathbf{o}} \ \mathbf{x}_t \ + \
  \mathbf{b}_{\mathbf{o}}
  &&&\text{output gate} 
\end{align}

mLSTM accommodates several memory cells similar to the standard LSTM architecture. In mLSTM, the presence of multiple heads or multiple cells yields the same outcome because memory mixing does not occur. To mitigate instability caused by the exponential gates in mLSTM, we apply the stabilization methods previously described for sLSTM (refer to Equation~\ref{eq:slstmstabil}). The absence of memory mixing in mLSTM enables the recurrence to be expressed in a parallelizable form. Additional information can be found in Appendix~\ref{sec:appmLSTM}.

\subsection{xLSTM Architecture}
\label{sec:xLSTMarch}

\paragraph{xLSTM Blocks.}
\label{sec:xLSTMblock}
The primary role of an xLSTM block is to encode preceding information through nonlinear transformations in a high-dimensional space, thereby enhancing the distinction between different context histories. The ability to disentangle various histories is fundamental for accurate prediction of subsequent sequence components, such as the next token. This approach leverages Cover’s Theorem~\citep{Cover:65}, which indicates that the chance of achieving linear separability among nonlinearly mapped patterns is increased in higher-dimensional representations as compared to the input space. We examine two designs for the residual block: (i) One design aligns with post up-projection used in Transformers, where the nonlinear summary of past information takes place in the original space, is then followed by a linear expansion into a higher-dimensional space, the application of a nonlinear activation, and finally, a linear reduction back to the original dimension; see the left portion of Figure~\ref{fig:Backbones} and the third column in Figure~\ref{fig:xlstm_sketch}. A more granular illustration appears in Figure~\ref{fig:appsLSTM_detailed} in the appendix. (ii) The other design involves a residual block featuring pre up-projection, similar to architectures found in State Space Models, where the input is first projected linearly into a high-dimensional space,

captures past information in a non-linear fashion within a high-dimensional representation, subsequently applying a linear transformation to revert to the original space. If an xLSTM block includes an sLSTM, the up-projection occurs after the core computation. Conversely, when utilizing an mLSTM within an xLSTM block, the up-projection precedes the primary module, as operating in a high-dimensional space enhances the block’s memory capacity. Additional visualizations can be found in the leftmost panel of Figure~\ref{fig:Backbones}, the third column of Figure~\ref{fig:xlstm_sketch}, and in Figure~\ref{fig:appsLSTM_detailed} located in the appendix.

\paragraph{xLSTM Architecture. }
\label{sec:xLSTMarchitecure}
The xLSTM architecture is formed by stacking building blocks in a residual manner~\citep{Srivastava:15,He:16}. Following prevailing approaches, we employ a residual backbone that incorporates pre-LayerNorm~\citep{Ba:16b}, as typically adopted in modern Large Language Models. Reference is made to the final column in Figure~\ref{fig:xlstm_sketch} for a visual summary.

\subsection{Memory and Speed Considerations}

Unlike Transformers, xLSTM networks feature linear computational complexity and fixed memory usage relative to sequence length. The compressive nature of xLSTM memory makes it particularly advantageous for real-world industry scenarios and for deployment on edge devices.

mLSTM utilizes memory that does not require trainable parameters, yet it incurs a high computational load owing to its $d \times d$ matrix memory and $d \times d$ update operations. This design balances an increase in memory capacity with higher computational demands. However, since these operations are highly amenable to parallelization on GPUs, the impact on overall runtime remains minimal.

mLSTM, much like FlashAttention~\citep{Dao:22,Dao:23} and GLA~\citep{Yang:23arxiv}, can be efficiently parallelized. In contrast, sLSTM cannot be parallelized as straightforwardly due to the hidden-hidden memory mixing. To address this, we introduced an optimized CUDA implementation leveraging GPU register memory, which achieves performance typically within a factor of two compared to mLSTM.

\section{Related Work}
\label{sec:related}

\paragraph{Linear Attention.} To address the issue of the Transformer's quadratic complexity with respect to context length, various approaches have been developed that achieve linear scaling of attention. The Synthesizer generates synthetic attention weights, foregoing token-to-token dependencies~\citep{Tay:20arxiv}. Linformer employs low-rank matrices to approximate self-attention in a linear manner~\citep{Wang:20arxiv}. Linear Transformer reforms the attention mechanism to make it inherently linear~\citep{Katharopoulos:20}. Performer leverages positive orthogonal random features to efficiently approximate the softmax operation within attention~\citep{Choromanski:21}. Alternative strategies, such as fast long convolutions, replace the traditional attention component in models like Structured Global Convolution (SGConv)~\citep{Li:22} and the Hyena Hierarchy~\citep{Poli:23}.

\paragraph{State Space Models.} State Space Models (SSMs) have garnered significant attention in recent times, owing to their linear computational complexity with respect to context length and competitive results in comparison to Transformer architectures. The sequence of notable developments in this domain began with the Structured State Space sequence model (S4)~\citep{Gu:21}, and was subsequently advanced by methods such as the Diagonal State Space (DSS) model~\citep{Gupta:22}, Gated State Space (GSS) models~\citep{Mehta:22}, S5 model~\citep{Smith:22}, Bidirectional Gated SSM (BiGS)~\citep{Wang:22}, H3 model~\citep{Fu:23}, and Mamba~\citep{Gu:24arxiv}.

\paragraph{Recurrent Neural Networks.} Recurrent Neural Networks (RNNs) have emerged as alternatives to Transformer and attention mechanisms, largely due to their linear scalability with sequence length. For language modeling tasks, RNNs equipped with Deep Linear Recurrent Units (LRUs) have demonstrated strong performance~\citep{Orvieto:23, De:24arxiv}. Hierarchically Gated Linear RNN (HGRN)~\citep{Qin:23} as well as its successor HGRN2~\citep{Qin:24arxiv} have also achieved encouraging results. Among RNN-based strategies for large language models, RWKV~\citep{Peng:23arxivshort,Peng:24arxivshort} stands out for its comparable effectiveness relative to Transformer architectures.

\paragraph{Gating.}
A central innovation in LSTM architectures is the gating mechanism, an idea that has been revisited and given new perspectives in numerous contemporary works. Gating has found applications in a variety of models, including HGRN~\citep{Qin:23}, HGRN2~\citep{Qin:24arxiv}, Gated Linear Attention (GLA)~\citep{Yang:23arxiv}, Gated State Space (GSS) models~\citep{Mehta:22}, Bidirectional Gated SSM (BiGS)~\citep{Wang:22}, Moving Average Equipped Gated Attention (MEGA)~\citep{Ma:22}, RWKV~\citep{Peng:23arxivshort}, and Mamba~\citep{Gu:24arxiv}.

\paragraph{Covariance Update Rule.}  
In order to improve storage capability, we have augmented the mLSTM cell by integrating a matrix memory using a covariance-based update mechanism. This principle also underpins several other approaches, such as Fast Weight Programmers~\citep{Schmidhuber:92ncfastweights,Schlag:21}, RWKV-5 and RWKV-6~\citep{Peng:24arxivshort}, Retention~\citep{Sun:23arxiv}, Linear Transformer~\citep{Katharopoulos:20}, and HGRN2~\citep{Qin:24arxiv}.

\paragraph{Most Related.} The models most similar in design to xLSTM are Retention~\citep{Sun:23arxiv}, RWKV~\citep{Peng:23arxivshort,Peng:24arxivshort}, and HGRN2~\citep{Qin:24arxiv}. All of these architectures implement matrix memory and/or gating mechanisms. Yet, unlike the novel sLSTM, they lack the capability for memory mixing. Memory mixing is critical for effectively addressing state tracking tasks, which renders LSTMs more powerful than State Space Models (SSMs) and Transformers~\citep{Merrill:24,Deletang:23}. Accurately tracking state is a prerequisite for tasks like code evaluation or maintaining coherence of entities in lengthy narratives.

\paragraph{Residually Stacking Architectures.} The xLSTM architecture, mirroring the design strategy of most modern large-scale deep learning systems, is built by residually stacking fundamental modules~\citep{Srivastava:15,He:16}.

This paradigm has been crucial to the success of deep convolutional neural networks~\citep{He:16} and, notably, Transformers~\citep{Vaswani:17}. Transformer-based models have become the primary foundation for Large Language Models (LLMs) such as GPT-3~\citep{Brown:20short}, ChatGPT~\citep{Schulman:22short}, GPT-4~\citep{Achiam:23gpt4short}, Megatron-LM~\citep{Shoeybi:19}, Gopher~\citep{Rae:21short}, ERNIE 3.0 Titan~\citep{Wang:21arxivshort}, GLaM~\citep{Du:21glamshort}, Chinese M6~\citep{Lin:21}, multilingual AlexaTM 20B~\citep{Soltan:22}, OPT~\citep{Zhang:22opt}, Chinchilla~\citep{Hoffmann:22short}, BLOOM~\citep{Scao:22arxivshort}, GLM-130B~\citep{Zeng:22arxivshort}, LaMDA~\citep{Thoppilan:22short}, PaLM~\citep{Chowdhery:22short}, Llama~\citep{Touvron:23arxiv}, and Gemini~\citep{GeminiTeam:23arxiv,Reid:24arxiv}.

\section{Experiments}
\label{sec:Exp}

This section presents an empirical assessment of xLSTM and benchmarks it against prevailing approaches in the domain of language modeling. Section~\ref{sec:ExpSyntethic} examines xLSTM's proficiency using designed synthetic challenges. 
Section~\ref{sec:ExpComparison} provides a comparison of validation perplexity for several leading language models, each trained on a 15B token corpus from SlimPajama~\citep{Soboleva:23}. Ablation experiments are conducted on xLSTM under identical conditions. Furthermore, we study how these models scale with increased data, drawing parallels to the methodologies in \citet{Kaplan:20} and \citet{Brown:20short}.
In Section~\ref{sec:ExpLanguage}, 
a comprehensive suite of language modeling tests is undertaken. Here, xLSTM is contrasted with the highest-performing models from Section~\ref{sec:ExpComparison}, following training on an expanded dataset of 300B tokens sourced from SlimPajama~\citep{Soboleva:23}. The evaluation includes: (1) analyzing the ability of models to generalize to longer contextual sequences, (2) comparing them based on validation perplexity and downstream task efficacy as described in~\citep{Sutawika:24}, (3) assessing performance across 571 distinct textual categories from the PALOMA language benchmark~\citep{Magnusson:23arxivshort}, and (4) revisiting scaling properties with the increased training data size, now totaling 20 times the original volume.

\label{sec:xlstm_blocks_notation}
Throughout all experiments, we denote xLSTM[$a$:$b$] as representing a ratio of $a/b$ where $a$ indicates the number of mLSTM-based blocks and $b$ represents the count of sLSTM-based blocks within the xLSTM architecture. As an illustration, xLSTM[7:1] specifies a configuration in which seven out of eight blocks are constructed with mLSTM, with the remaining block composed of sLSTM. In the commonly used setup involving a total of 48 blocks, this ratio yields 42 mLSTM-based and 6 sLSTM-based blocks. Additionally, across all experiments, we incorporate up-projection blocks before and after mLSTM and sLSTM blocks, respectively.

\subsection{Synthetic Tasks and Long Range Arena}
\label{sec:ExpSyntethic}
Initially, we evaluate the capability of xLSTM's recently introduced exponential gating combined with memory mixing by applying it to formal language tasks~\citep{Deletang:23}. Subsequently, we investigate how well xLSTM's innovative matrix memory mechanism performs on the Multi-Query Associative Recall benchmark~\citep{Arora:23arxiv}. Lastly, we measure xLSTM's aptitude for handling lengthy sequences by benchmarking it on the Long Range Arena suite~\citep{Tay:21}.

\paragraph{Evaluation of xLSTM's Exponential Gating with Memory Mixing.} We examine the capabilities of xLSTM's novel exponential gating approach combined with memory mixing, which is anticipated to facilitate effective state tracking~\citep{Merrill:24, Merrill:23}. To this end, we adapt and expand the suite of formal language benchmarks established by \citet{Deletang:23}, allowing for training across multiple sequence lengths, thereby supporting length generalization studies. Comprehensive task specifications and supplementary results are provided in Appendix~\ref{sec:appExpSynthetic-formalLang}. 
Our evaluation benchmarks xLSTM against alternative architectures, including Transformers, State Space Models, and traditional Recurrent Neural Networks. For each method, performance is quantified as accuracy on task-relevant tokens, with scores normalized such that 0 corresponds to chance level and 1 to perfect accuracy. Specifically, we investigate 2-block variants of the following models: xLSTM[0:1] (pure sLSTM), xLSTM[1:0] (pure mLSTM), xLSTM[1:1], Llama, Mamba, RWKV, Retention, Hyena, LSTM, and LSTM implemented within Transformer-style blocks (LSTM (Block)). The comparative findings are presented in Figure~\ref{fig:formal-main}. 
It is observed that architectures lacking memory mixing—such as standard Transformers or State Space Models—are unable to address tasks requiring state tracking, such as regular grammar tasks exemplified by parity detection. This observation aligns with previous reports that highlight the inherent limitations of Transformers and State Space models relative to RNNs in terms of computational expressivity~\citep{Merrill:24, Merrill:23, Deletang:23}.

\paragraph{Test of xLSTM's Memory Capacities on Associative Recall Tasks.} In this study, we evaluate the matrix-based memory mechanism introduced in xLSTM by assessing its memory capacity on the Multi-Query Associative Recall task~\citep{Arora:23arxiv}. This benchmark requires models to memorize randomly selected key-value pairs from a sizable vocabulary within each sequence for subsequent retrieval. To further raise the challenge beyond the original design, we scale the number of key-value pairs up to 256 and extend the possible context length to 2048, thereby facilitating a more comprehensive assessment of the memory limits across various models. Specifically, we examine two-block versions of Llama, Mamba, RWKV-5, RWKV-6, xLSTM[1:1], and xLSTM[1:0]. Performance is quantified by measuring the accuracy with which models recall these associations. Since Transformer-based models such as Llama possess memory scaling exponentially with the representation dimension~\citep{Ramsauer:21}, they serve as a reference point in this evaluation. Figure~\ref{fig:mqar-main} illustrates the results.
Among models outside of the Transformer family, xLSTM[1:1] attains the highest recall performance, even for models with fewer parameters. Notably, integrating the sLSTM block does not curtail memory capacity; in fact, it amplifies it, as demonstrated in the most demanding scenario involving 256 key-value pairs. Supplementary findings can be found in Appendix~\ref{sec:appExpSynthetic-mqar}, where extrapolation studies reveal that xLSTM's advanced memory capacities support generalization to longer context lengths beyond those encountered in training.

\paragraph{Test of xLSTM's Long Context Capabilities on Long Range Arena.} To evaluate how well xLSTM processes extensive sequences and broad contextual information, we benchmark it against other models using the Long Range Arena~\citep{Tay:21}. Across the various tasks, xLSTM reliably achieves high performance, indicating that the architecture is particularly effective in managing diverse long-context challenges.

For more details, see Appendix~\ref{sec:appExpSynthetic-lra}.

\subsection{Method Comparison and Ablation Study}
\label{sec:ExpComparison}

In order to investigate our central research question—specifically, the capabilities of our proposed LSTM variants when scaled up for language modeling—we conduct experiments by training xLSTMs, Transformers, State Space Models, and additional architectures on 15B tokens from SlimPajama, all within an auto-regressive framework. Following training, we evaluate model performance on the validation dataset and carry out ablation experiments with xLSTMs.

\paragraph{Evaluation of xLSTM Against Existing Approaches.} To benchmark performance, models are trained on a dataset consisting of 15 billion tokens from SlimPajama~\citep{Soboleva:23}, and their validation perplexities are measured. The following models are included in the analysis: xLSTM (our proposed approach), GPT-3 (Transformer architecture)~\citep{Brown:20short}, Llama (Transformer)~\citep{Touvron:23arxiv}, H3 (State Space Model)~\citep{Fu:23}, Mamba (State Space Model)~\citep{Gu:24arxiv}, RWKV-4 (RNN)~\citep{Peng:23arxivshort}, RWKV-5 (RNN)~\citep{Peng:24arxivshort}, RWKV-6 (RNN)~\citep{Peng:24arxivshort}, GLA (linear Transformer)~\citep{Yang:23arxiv}, HGRN (RNN)~\citep{Qin:23}, HGRN2 (RNN)~\citep{Qin:24arxiv}, RetNet (linear Transformer)~\citep{Sun:23arxiv}, Hyena (linear Transformer)~\citep{Poli:23}, xLSTM[1:0], and xLSTM[7:1] (refer to Section~\ref{sec:xlstm_blocks_notation} for block notation details). All models employ mixed precision training. However, RWKV-5, RWKV-6, GLA, and HGRN2 use an alternative mixed precision implementation, not relying on PyTorch’s automatic mixed precision mechanism (see Appendix Section~\ref{sec:appGenTrainProc}). 

Approaches are grouped as follows: (a) Transformers, (b) State Space Models (SSMs), and (c) both Recurrent Neural Networks (RNNs) and linear Transformer methods. The linear Transformer variants replace the attention operation with linear alternatives. Each architecture is scaled to match a GPT-3 configuration with 350 million parameters, corresponding to an embedding dimension of 1024 and 24 residual blocks. Notably, only GPT-3 shares token and output embeddings, resulting in a reduced parameter count compared to others. According to the outcomes in Table~\ref{tab:spaj15b_model_results}, xLSTM achieves the lowest validation perplexity, surpassing all other compared models. For further information, refer to Appendix~\ref{sec:appExpComparison}. Figure~\ref{fig:temp_sclaw15B} visualizes model scaling trends, suggesting that xLSTM maintains its superior performance as model size increases.

\paragraph{Ablation Studies.}
Table~\ref{tab:spaj15b_model_results} and Figure~\ref{fig:temp_sclaw15B} illustrate that xLSTM delivers strong performance on language modeling tasks after being trained with 15B tokens sourced from SlimPajama. To investigate the impact of each architectural modification from LSTM to xLSTM, we systematically convert a standard LSTM into an xLSTM. The initial step involves embedding LSTM layers within pre-LayerNorm residual backbone structures. Next, we augment the architecture by implementing a post up-projection block. Lastly, exponential gating and matrix memory components are incorporated.

The results are shown in Table~\ref{tab:ablstudies} (top). 

Ablation experiments reveal that the significant enhancement in performance can be credited to the interplay between exponential gating and the matrix memory. Furthermore, recognizing the crucial role that gating plays in both RNNs and State Space Models, we systematically evaluate several gating strategies. The results summarized in Table~\ref{tab:ablstudies} (bottom) indicate that when gates are individually learnable and conditioned on the input, there is a progressive benefit in model performance. Further investigations concerning specific backbone elements are presented in Appendix~\ref{sec:appAblDetails}.

\subsection{xLSTM as Large Language Model}
\label{sec:ExpLanguage}

The final phase of our investigation involves conducting extensive language modeling experiments to evaluate xLSTM as a large language model (LLM). For this purpose, we scale up the training process, utilizing 300B tokens sourced from SlimPajama. This token count aligns with that employed by models such as Mamba~\citep{Gu:24arxiv} and Griffin~\citep{De:24arxiv}.

Our comparative analysis includes xLSTM alongside RWKV-4, Llama, and Mamba, with each model chosen to represent their respective categories discussed in Section~\ref{sec:ExpComparison}. RWKV-4 is designated as the RNN benchmark due to recent advancements in training precision for RWKV-5, RWKV-6, and HGRN2, which were established only after the commencement of the 300B token trials (refer to Appendix~\ref{sec:appGenTrainProc}).

We examine multiple model scales (125M, 350M, 760M, 1.3B), evaluate all configurations for their ability to extrapolate sequence length, and measure validation set accuracy. The models' effectiveness is further assessed through downstream task evaluation, language modeling tests across 571 distinct domains from the PALOMA benchmark, and an analysis of scaling law characteristics.

\paragraph{Sequence Length Extrapolation.}
\label{sec:spaj300B_ctx_extrapolation}
To begin, we assess the sequence length extrapolation capabilities of large-scale (1.3B parameter) models, specifically xLSTM, RWKV-4, Llama, and Mamba. Each model undergoes training with a fixed context window of 2048 tokens, after which we evaluate them on sequences extending to 16384 tokens. The findings are visualized in Figure~\ref{fig:spaj300B_seq_extrapolation}. Unlike the other architectures, xLSTM preserves lower perplexity scores even as the input sequence length increases substantially.

\paragraph{Validation Perplexity and Downstream Tasks.}
\label{sec:spaj_downstream_eval}
Next, across all model capacities, we conduct an evaluation of xLSTM, RWKV-4, Llama, and Mamba on the SlimPajama validation dataset to gauge their next token prediction accuracy, as well as their performance on additional tasks targeting common sense reasoning.

The third column in 
Table~\ref{tab:lmeval} presents the perplexity scores on the validation set for various approaches. Both xLSTM[1:0] and xLSTM[7:1] achieve the lowest perplexity across all model sizes, signifying their superiority in this regard. Performance metrics for downstream tasks are reported in the remaining columns of 
Table~\ref{tab:lmeval}. Across nearly all tasks and model scales, xLSTM consistently outperforms competing methods; exceptions are found only in the ARC task, where Mamba sometimes surpasses xLSTM. Further information can be found in Appendix~\ref{sec:appExpLanguage}.

\paragraph{Performance on PALOMA Language Tasks.} Next, we evaluate how xLSTM, RWKV-4, Llama, and Mamba perform at next token prediction across all model sizes using the PALOMA language tasks~\citep{Magnusson:23arxivshort}. We assess each model's effectiveness via perplexity scores on 571 distinct text domains, ranging from nytimes.com to r/depression on Reddit. 

Table~\ref{tab:ppleval} displays perplexity results, divided into categories for language modeling (first seven columns) and specific domain benchmarks (final five columns). On these language tasks, xLSTM[1:0] consistently achieves lower perplexity compared to xLSTM[7:1]. 

Specifically, xLSTM[1:0] surpasses Mamba in perplexity across 568 of 571 (99.5\%) domains, outperforms Llama on 486 of 571 (85.1\%), and exceeds RWKV-4 on 570 of 571 (99.8\%) domains. A more comprehensive breakdown can be found in Appendix~\ref{sec:appExpLanguage}.

\paragraph{Scaling Laws.} In our fourth analysis, we investigate the power-law scaling dynamics, which permit projecting model performance as architectures expand in size~\citep{Kaplan:20,Brown:20short}. The scaling curves depicted in Figure~\ref{fig:temp_sclaw300B} reveal that while all evaluated models exhibit comparable scaling trends, they differ in their baseline performance. Among these, RWKV-4 demonstrates the lowest performance, followed sequentially by Llama and then Mamba. The xLSTM architecture surpasses Mamba, with the performance margin between xLSTM and Mamba roughly equal to that observed between Mamba and Llama. These scaling trends suggest that as models grow in parameter count, xLSTM maintains a relative advantage over both Transformer and State-Space approaches.

\textbf{Generation Times and Maximal Throughput.} In Figure~\ref{fig:inference_time_generation}, we evaluate the generation times for our xLSTM models at the 1.3B parameter level, and in Figure~\ref{fig:inference_time_decoding_throughput} (left), we assess their maximal throughput. For reference, we include equivalently sized Mamba, Llama, and RWKV models using implementations from HuggingFace; in particular, for Llama we incorporate the static key-value cache. During our testing, it was not feasible to use fully compiled key-value caches for the Transformer Model, nor was Mamba compatible with \texttt{torch.compile}. All models in our text generation benchmarks were evaluated with a batch size of 1 and a pre-fill of 16, a configuration that provides an optimal advantage for the Transformer architecture.

The results in Figure~\ref{fig:inference_time_generation} make evident that xLSTM, as well as the other recurrent architectures Mamba and RWKV-4, exhibit generation times that scale linearly, whereas Llama demonstrates the expected quadratic increase. In assessing decoding throughput, we varied batch sizes and prefill values specifically for the Llama model. Figure~\ref{fig:inference_time_decoding_throughput} (right) further demonstrates that xLSTM supports much larger batch sizes than Llama, benefiting from constant memory requirements and thereby yielding the highest throughput among the compared models.

\section{Limitations}

(i) Unlike mLSTM, the sLSTM’s approach to memory mixing necessitates sequential processing, which restricts the possibility of parallel execution and thus limits efficient parallelization. Despite this, we implemented an expedited CUDA kernel for sLSTM that operates at less than double the time required by our parallelized mLSTM version. (ii) The current CUDA kernels for mLSTM lack advanced optimization, resulting in runtimes that are approximately four times slower compared to FlashAttention or the scan algorithm utilized in Mamba. Enhanced CUDA kernels modeled after FlashAttention could potentially improve these speeds. (iii) The matrix memory of mLSTM entails high computational cost because processing involves $d \times d$ matrices. However, since memory updates and accesses are parameter-free and utilize standard parallel matrix operations, the additional runtime burden due to complex memory structure remains minor. (iv) Proper selection of initialization values for the forget gates is critical. (v) Because the matrix memory is decoupled from sequence length, extending the sequence can strain memory capacity in cases of longer contexts. Nonetheless, for context sizes up to 16k, this issue is not significant, as discussed in Section~\ref{sec:spaj300B_ctx_extrapolation}. (vi) Due to resource demands in large language model training, neither the architecture nor the hyperparameters were extensively tuned, particularly for larger xLSTM configurations. Substantial optimization efforts will be necessary for xLSTM to realize its maximum performance.

\section{Conclusion}
We have provided a partial response to our primary inquiry: What progress can be achieved in language modeling by expanding LSTM architectures to encompass billions of parameters? The evidence thus far demonstrates that this approach yields results comparable to those obtained with contemporary architectures such as Transformers and State Space Models. Our contributions include augmenting LSTM into xLSTM, incorporating exponential gating mechanisms with memory mixing and introducing an innovative memory structure. Comparative assessments show that xLSTM models excel in language modeling tasks relative to prevailing approaches like Transformers and State Space Models. Observed scaling laws suggest that with increased model size, xLSTM emerges as a formidable challenger to the dominant Large Language Models utilizing Transformer frameworks. Furthermore, xLSTM holds promise for transformative advances in other areas, including Reinforcement Learning, Time Series Prediction, and the modeling of physical systems.

\end{document}